\chapter{System Design and Methodology}
\label{chap:architecture}

\section{System Architectural Overview}
\label{sec:system_overview}

Olit is designed with a clear separation between the user interface and the AI's decision-making logic. Figure \ref{fig:arch_diagram} outlines the system architecture, illustrating the strict boundary between the local browser environment and the remote Galaxy server.

\begin{figure}[htbp]
\centering
\fbox{\parbox{0.9\textwidth}{\centering
\vspace{0.5cm}
\textbf{[Client Browser Environment]}\\[0.2cm]
\fbox{UI Shell: Orbit Chat Panel + Artifact Pane + Modals}\\[0.3cm]
$\Updownarrow$ \textit{(Asynchronous Worker Messaging)}\\[0.3cm]
\fbox{Agent Engine: pi-durable Web Worker}\\[0.2cm]
$\Downarrow$ \hspace{2cm} $\Downarrow$ \hspace{2cm} $\Downarrow$\\[0.2cm]
\fbox{Pyodide Wasm} \quad \fbox{OPFS SQLite} \quad \fbox{Skills Router}\\[0.5cm]
\hrulefill\\[0.3cm]
\textbf{[Remote Galaxy Server Infrastructure]}\\[0.2cm]
$\Uparrow$ \textit{(Authenticated REST APIs using Session Cookies)}\\[0.2cm]
\fbox{Galaxy API: Tools, Datasets, Workflows, Pages, Charts}
\vspace{0.5cm}
}}
\caption{Architectural decomposition of Olit, demonstrating client-side execution over remote Galaxy server infrastructure.}
\label{fig:arch_diagram}
\end{figure}

The codebase comprises two primary subsystems:
\begin{enumerate}
    \item \textbf{The UI Shell} 
    \item \textbf{The Agent Engine}
\end{enumerate}

\section{Human-in-the-Loop}
\label{sec:governance_plan_first}

A critical vulnerability of naive language model agents is precipitous action: an unconstrained agent given high-level goals may immediately execute unverified tools, consume substantial cluster allocations, or inadvertently overwrite research data.
\subsection{The Execution Lifecycle}
To prevent unintended actions, the agent is restricted from executing tools immediately after receiving a user request. Instead, it follows a structured lifecycle:
\begin{enumerate}
\item \textbf{Data Exploration:} The agent first inspects dataset metadata (such as column names, datatypes, and sample counts) using read-only operations.
\item \textbf{Plan Formulation:} It then drafts a step-by-step execution plan. This includes specific tool choices (e.g., FastQC, Trimmomatic) and a detailed table of all proposed parameters.
\item \textbf{User Review:} The proposed plan is displayed in the UI, giving the user the option to \textbf{Approve}, \textbf{Edit}, or \textbf{Reject} it. The agent pauses completely until the user makes a decision.
\item \textbf{Execution and Tracking:} Once approved, the agent submits the tools to the Galaxy API in the correct order and continuously monitors their progress.
\end{enumerate}

\subsection{Safety Gates for Destructive Actions}
Operations that permanently delete data---such as purging histories, deleting datasets, or stopping active workflows---are handled by a dedicated safety mechanism (\texttt{destructiveGate}). Even if the agent is already in the middle of an approved execution phase, attempting a destructive action automatically triggers a confirmation modal. If the user declines, the action is blocked to prevent any accidental data loss.

\section{The Dynamic Skills Knowledge Architecture}
\label{sec:skills_system}

\subsection{The On-Demand Router Pattern}
Olit implements an on-demand \textbf{Skills Router} (\texttt{skills.ts}). 

By loading procedural rules only when active, the agent maintains maximum attentional focus on the immediate operational task.

\subsection{Skills Repositories and Cryptographic Vendoring}
Olit organizes procedural knowledge across two repositories:
\begin{itemize}
    \item \textbf{\texttt{olit-skills}:} 
    \item \textbf{\texttt{galaxy-skills}:} 
\end{itemize}


\section{Data Visualization and Interactive Analytics}
\label{sec:visual_analytics}

Data exploration in Olit is managed by a dual-pathway architecture implemented in \texttt{visualizations.ts} and \texttt{vega.ts}


\section{Session Resumability and Automated Notebooks}
\label{sec:resumability_notebooks}

\begin{figure}[htbp]
\centering
\fbox{\parbox{0.9\textwidth}{\centering
\vspace{0.4cm}
\textbf{[User Opens Olit on History]}\\[0.2cm]
$\Downarrow$\\[0.2cm]
\fbox{Local Browser Storage (OPFS SQLite) contains session?}
\begin{tabular}{cc}
\textbf{YES} $\longrightarrow$ & \fbox{Restore Full Conversation Transcript Immediately} \\
\textbf{NO} $\longrightarrow$ & \fbox{Check Galaxy Server for Attached History Notebook}
\end{tabular}\\[0.2cm]
$\Downarrow$\\[0.2cm]
\fbox{Existing Notebook Found in History?}
\begin{tabular}{cc}
\textbf{YES} $\longrightarrow$ & \fbox{Attach to Existing Notebook, Resume Analysis State} \\
\textbf{NO} $\longrightarrow$ & \fbox{Initialize New Conversation and Create Fresh Notebook}
\end{tabular}\\[0.2cm]
$\Downarrow$\\[0.2cm]
\fbox{Active Work and Iterative Analysis}\\[0.2cm]
$\Downarrow$\\[0.2cm]
\fbox{User Saves Session?}
\begin{tabular}{cc}
\textbf{YES} $\longrightarrow$ & \fbox{Serialize Session to Galaxy Visualization Object (Shareable)} \\
\textbf{NO} $\longrightarrow$ & \fbox{Session Persists Securely in Local Browser Cache}
\end{tabular}
\vspace{0.4cm}
}}
\caption{Decision tree for Olit session resumption across browser tabs, local storage, and Galaxy server instances.}
\label{fig:resumability_flow}
\end{figure}

If the user revisits the analysis within the same browser, the complete conversational transcript is immediately restored from local SQLite storage. If accessed from an alternative machine or a fresh browser profile, Olit checks for an attached Galaxy Page notebook; if found, it parses the completed analytical history and resumes seamlessly.

\subsection{Automated Laboratory Notebook Generation}
Chat interfaces are not well-suited for long-term scientific record-keeping. To ensure the analysis is permanently documented, Olit automatically generates and updates a \textbf{Galaxy Page} in the background:
\begin{itemize}
    \item \textbf{Execution Tracking:} The approved plan is documented as a live Markdown checklist. As jobs finish, the agent updates the status of each step, marking them as complete (\texttt{- [x]}) or failed (\texttt{- [!]}).
    \item \textbf{Persistent Dataset Linking:} To guarantee exact reproducibility, datasets are referenced using both their simple UI numbers (HIDs) and their permanent, unique internal Galaxy identifiers.
    \item \textbf{Embedding Visualizations:} Interactive charts and figures generated during the session are directly embedded into the page using \texttt{{{artifact}}} tags. This turns the session into a clean, shareable summary of the entire workflow.
\end{itemize}

